\chapter*{Abstract}
\addcontentsline{toc}{chapter}{Abstract}

Dark patterns are deceptive or manipulative interface designs that steer users toward decisions they may not have otherwise intended to make. Prior research has shown that dark patterns can be systematically identified, categorized, and measured across digital interfaces, particularly in web and consent-related settings \cite{gray2018darkpatterns,mathur2019darkpatterns,nouwens2020gdpr}. However, most existing studies focus on whether dark patterns exist and how they should be classified, while paying comparatively less attention to whether different user profiles may experience different levels or types of dark pattern exposure.

This report addresses this gap by proposing a profile-aware auditing framework for measuring dark pattern exposure in mobile applications as a trajectory-based, profile-conditioned outcome. Rather than relying on demographic attributes, the framework constructs controlled profiles through two reproducible components: interest distributions over profile-specific keywords and behavior distributions over interaction actions. These profiles generate comparable app-interaction trajectories while keeping the audit focused on observable behavior and interest conditions.

For each recorded screen, the audit assigns a seven-dimensional DPS vector covering Interface Interference, Friction, Pressure, Contextual Deception, Sneaking, Nagging, and Forced Action. The empirical study covers 16 apps, 3 profiles, and 30 steps for each app-profile pair, producing 1,440 interaction records and 10,080 DPS labels. The analysis estimates profile-conditioned exposure \(E(D \mid A)\), compares differential exposure between profiles, and evaluates observed gaps using heatmaps, running mean convergence, permutation tests, corrected q-values, L1 distance, and Cramér's V.

The results suggest that dark pattern exposure is not uniformly distributed across controlled profiles. The high-payment shopping and bargain-hunter profiles tend to encounter stronger Pressure and Interface Interference exposure, while low-payment browsing still receives meaningful exposure in dimensions such as Forced Action and Nagging. Exposure estimates also become more stable after repeated interaction steps. These findings should be interpreted as evidence of profile-conditioned exposure differences in the observed audit data, not as proof of causality or internal platform intent. The contribution of this report is a reproducible profile-aware auditing method that extends dark pattern research from presence-oriented detection toward comparative measurement of profile-conditioned and differential exposure. The source code and experimental results are available at \url{https://github.com/santal0/OSN-PDEA}.
